%!TEX root = ../template.tex %
% !TeX spellcheck = en_US
%%%%%%%%%%%%%%%%%%%%%%%%%%%%%%%%%%%%%%%%%%%%%%%%%%%%%%%%%%%%%%%%%%%%
%% abstract-en.tex
%% NOVA thesis document file
%%
%% Abstract in English
%%%%%%%%%%%%%%%%%%%%%%%%%%%%%%%%%%%%%%%%%%%%%%%%%%%%%%%%%%%%%%%%%%%%

\typeout{NT FILE abstract-en.tex}%

In this thesis, we use a combination of analytical and numerical methods to evaluate the excitonic and optical effects in isolated or nanostructured monolayered semiconductors or insulators.
Systems in two dimensions capable of presenting many-body effects, through strong optical responses and highly dispersive surface-confined light--matter hybrids, caused a shift in the attention of the scientific community to
two-dimensional systems, and how to tailor them for applications in nanophotonics, polaritonics, optoelectronics, optical sensors, among others.
The capacity of two-dimensional systems in supporting long-lived excitons at room temperatures makes them viable for fabricating devices that could operate outside cryogenic conditions.

This thesis proposes ways of manipulating material properties by engineering the experimental setup, providing means of predicting the optical properties of two-dimensional materials with a bandgap more efficiently.
For that, we guide the reader through the theoretical concepts that we need to move forward the ideas herein explored.
The starting point is the classical formulation of electromagnetic phenomena in solids, including the application of the formalism to two-dimensional materials.
Next, we introduce the formalism of quantum mechanics to describe the physics of solid matter at the microscopic level, to provide insights the classical formalism does not offer, and to prepare the reader for the rest of the dissertation.
In particular, we provide the necessary background to describe the concept of the exciton, which is of extreme importance in this work.
Afterward, the theory behind the excitons is formulated, making it possible to study optical effects in two-dimensional materials caused by excitonic effects.

To choose the right material with practical applications in mind is not trivial.
Furthermore, it has been seen before in the literature that monolayer \gls{hbn} presents strong optical resonances in the ultraviolet.
We propose a way to engineer the optical properties of this material by setting up a nanostructured environment, where the monolayer is deposited in a substrate in contact with a metallic grating.
The idea is to change the electronic and optical properties of \gls{hbn} and how they interplay, by subjecting it to an electrostatic periodic potential.
In this way, we can change the bandgap as we like, enabling changing the excitonic and optical properties as well.
By tuning the potential, we can make \gls{hbn} optically useful in a wider range of frequencies down to the mid-infrared, increase the oscillation strength of the optical resonances, enhance its optical absorption,
and make it promising for application in nanophotonics and polaritonics.
All of this, using always the same material.

Finally, we look at the problem of electrostatic screening in single-layer semiconductors and insulators from a strict two-dimensional perspective.
Basically, we can reformulate the problem of calculating numerically the dielectric function the way it is done in the existing literature, by reducing the dimensionality of the approach, from three to two dimensions.
The theory applied to monolayer hBN produces results that agree satisfactorily with the ones reported in previous works, and the gain in numerical efficiency is notable.
Then, the problem is reformulated once more; we go from a strict two-dimensional description to a quasi-two-dimensional one.
By taking the thickness of the monolayer into account, we improve the agreement between our results and the ones we compared with previously.
This is allowed by a mixed-representation of the dielectric function, defining it in a combination of the two-dimensional reciprocal space with the one-dimensional space perpendicular to the layer.
The code developed is analyzed and validated, this time for a different material, monolayer molybdenum disulfide, \ch{MoS2}.
Once more, the results from our approach agree very well with more sophisticated ones, validating our theory.

